\documentclass[10pt]{article}
\usepackage[margin=0.75in]{geometry}
\usepackage{lmodern,amsmath,booktabs,array,enumitem,xcolor,hyperref}
\hypersetup{colorlinks=true,linkcolor=blue,urlcolor=blue}
\setlist[itemize,enumerate]{nosep,leftmargin=*}
\setlength{\parskip}{5pt}\setlength{\parindent}{0pt}
\title{Near-RT RIC and xApp Design for the\\Multi-Tier Digital Twin v3}
\author{Rao Yenamandra --- \texttt{raosy@digitaltwinsim.com}\\
Mubanga Nsofu --- \texttt{mubanga.nsofu@vodacom.co.za}\\
Asokan Ram --- \texttt{asokan.ram@wrc-nc.org}}
\date{Version 3 --- September 2026}
\begin{document}\maketitle

\section{Purpose and scope}
This document defines how the Digital Twin v3 maps its mIoT PRACH-overload experiment to an O-RAN deployment. The proposed PPO or Double-DQN controller is an \textbf{xApp} hosted by the \textbf{Near-Real-Time RAN Intelligent Controller (Near-RT RIC)}. It observes aggregated random-access measurements and recommends bounded overload-control parameters. The serving gNB/O-DU retains final enforcement authority.

This is a design mapping for a synthetic screening simulator. It does not claim that every named measurement or control is already exposed by a specific vendor E2 service model, nor that the browser model is a protocol-conformance implementation.

\section{Where the components are located}
\begin{tabular}{>{\raggedright\arraybackslash}p{0.22\textwidth}>{\raggedright\arraybackslash}p{0.30\textwidth}>{\raggedright\arraybackslash}p{0.40\textwidth}}
\toprule Component & Typical location & Responsibility in this design\\\midrule
UE / mIoT device & Customer or field device & Detects access opportunity, applies broadcast barring, selects a preamble, transmits and retries after failure/backoff.\\
O-RU & Cell site & Radio-frequency transmission/reception; no PPO/DQN decision.\\
O-DU / gNB & Cell site or far-edge cloud & Executes PRACH, detects preambles, produces RAO counters, applies approved ACB/backoff settings and preserves safety/fallback control.\\
Near-RT RIC & Operator edge or regional cloud, possibly near the O-CU & Hosts xApps, collects E2 telemetry, arbitrates policies and sends authorized control recommendations on a near-real-time timescale.\\
PRACH xApp & Software inside Near-RT RIC & Runs PPO, Double DQN, or the selected deterministic comparator; converts observations into bounded PRACH mitigation actions.\\
Non-RT RIC/rApp and SMO & Operator management/regional cloud & Offline analytics, policy, model training/validation, versioning and deployment approval.\\
\bottomrule\end{tabular}

\section{Closed-loop workflow}
\begin{enumerate}
\item The O-DU/gNB aggregates attempted preambles, occupied/idle preambles, collision estimate, successful accesses, failures, retries and waiting-load indicators for each cellular mIoT service area.
\item Measurements are delivered toward the Near-RT RIC through an E2-based telemetry path. The simulator represents this as one observation per control interval rather than packet-level E2 messages.
\item The xApp builds a state containing backlog, collision fraction, idle-preamble indication, storm phase, elapsed RAOs and current control settings.
\item The selected policy produces an action. PPO and Double DQN select from the same bounded action set; non-adaptive, demand-follow and rule-based controllers provide reproducible comparators.
\item A guard validates range, rate of change, freshness and policy authority. Unsafe, stale or unavailable output is rejected and the gNB retains its deterministic fallback.
\item The approved recommendation is sent to the serving gNB/O-DU. It applies the new setting at an allowed configuration boundary, never to an individual in-flight PRACH attempt.
\item Subsequent RAO measurements close the loop and are logged for evaluation and later retraining.
\end{enumerate}

\section{xApp observations, actions and safeguards}
\textbf{Observations} include RAO-normalized backlog, admitted attempts, collision fraction, idle-preamble fraction, successes, exhausted attempts, mean retry count, access delay, storm indicator and prior action. They are aggregated for the mIoT slice; eMBB/URLLC scheduled traffic and Wi-Fi offload are outside this PRACH loop.

\textbf{Actions} are bounded combinations of ACB admission probability, barring window and retry-backoff window. A future implementation may also expose standards/vendor-supported RAO or preamble configuration, but the v3 comparison does not assume instantaneous physical-resource reconfiguration. Demand-follow uses $p_{\mathrm{ACB}}=\operatorname{clip}(\alpha M/B,0.005,1)$, where $M$ is contention preambles, $B$ backlog and $\alpha$ the editable target load.

\textbf{Guards} enforce minimum/maximum values, maximum step change, minimum dwell time, telemetry freshness, rollback, audit logging and operator priority. Loss of the RIC/E2 path must not stop random access: the gNB falls back to its last approved or rule-based configuration.

\section{Training and inference split}
Training is performed offline or under an rApp/Non-RT RIC workflow using Digital Twin scenarios, multiple random seeds and held-out storms. Candidate models are checked against deterministic controllers before approval. The deployed xApp performs inference and limited bounded adaptation; it does not conduct unconstrained online exploration on a live RAN. Model identity, training data, reward weights, guard configuration and rollback version are recorded.

\section{Controller behaviour in the simulator}
\begin{tabular}{>{\raggedright\arraybackslash}p{0.20\textwidth}>{\raggedright\arraybackslash}p{0.16\textwidth}>{\raggedright\arraybackslash}p{0.56\textwidth}}
\toprule Controller & Adaptive online? & Function\\\midrule
Non-adaptive & No & Holds fixed slice allocation, ACB and backoff settings.\\
Demand-follow & Yes & Scheduled PRBs follow observed slice demand; separately, mIoT ACB follows PRACH backlog using a fixed formula.\\
Rule-based & Yes & Changes bounded controls when collision/load thresholds are crossed.\\
PPO xApp & Yes & Selects a bounded action from a trained stochastic policy using current state.\\
Double-DQN xApp & Yes & Selects the highest-valued bounded action using current state.\\
Offline optimizer & No after deployment & Searches for a better fixed configuration before operation.\\
\bottomrule\end{tabular}

\section{Reward, decision and reported KPIs}
The PRACH contribution rewards successful access and penalizes collisions, unresolved backlog and exhausted attempts. The overall research score also retains slice satisfaction and continuity terms. A learned controller is not declared superior merely because reward rises: the comparison reports access-success rate, failures, collision rate, retry transmissions, mean and 95th-percentile delay, peak backlog and clearance RAO against the best deterministic controller and the other learned controller.

\section{Boundaries with other v3 functions}
PRACH mitigation applies only to the cellular mIoT slice on mIoT-capable tiers. The cellular mmWave hotspot/small-cell tier carries eMBB/URLLC and is excluded. Wi-Fi offload is a separate macro/micro-to-Wi-Fi eMBB access-selection workflow using discovery, RSSI, authentication and load gates. The Near-RT PRACH xApp does not control autonomous UE Wi-Fi selection.

\section{Path from simulator to field trial}
Before field use, map each observation and action to supported E2 service-model/vendor interfaces; calibrate arrivals and collision inference with lab or carrier traces; define control periodicity and delay budgets; test E2 loss and rollback; run shadow-mode recommendations; then conduct a bounded canary trial. The v3 browser result alone is evidence of reproducible algorithm behaviour, not production readiness.

\end{document}
